구간 합은 어떤 배열의 $s$번째 인덱스 값부터 $e$번째 인덱스의 값을 빠르게 구하는 방법을 말한다. 만약 반복문을 이용해 전부 다 더한다면 $e-s$번의 연산이 필요하겠지만, 구간 합을 이용하면 $1$번의 연산으로 해결할 수 있다. 구간 합은 단순하지만 많은 PS 문제, 특히 동적계획법 문제에서 자주 사용되므로 꼭 숙지하도록 하자.

\section{부분 합을 구하는 문제}
다음 문제를 생각해보자.

\begin{example} \label{ex: 1d-sum-query}
    (11659번) 길이 $N$인 어떤 배열 $A$에서 다음 쿼리 $Q$개를 해결하는 프로그램을 작성하라.
    \begin{itemize}
        \item 쿼리: $A_s + A_{s+1} + \cdots + A_e$를 구하여라.
    \end{itemize}
\end{example}

만약 반복문을 통해 구한다면 $O(NQ)$의 시간복잡도를 가질 것이다. 그런데 잘 생각해보면 낭비되는 계산이 있다. 만약 $s=3, e=10$인 쿼리와 $s=4, e=11$인 쿼리를 처리해야 한다고 생각해보자. 그러면 $s=4$에서 $e=10$인 부분의 덧셈은 두 번 진행되고 이는 매우 불필요하다!

이 중복 문제를 해결하기 위해 누적 합 $S_i$를 생각해보자. 누적 합 $S_i$는 $S_i - S_{i-1} = A_i, S_0 = 0$으로 정의되는 수열이다. $S_i$의 의미는 $1$번째부터 $i$번째까지 수의 합이다. 이를 이용하면 $s$번째 수부터 $e$번째 수까지의 합을 $S_e - S_{s-1}$로 구할 수 있다. 구현은 다음 코드 \Cref{code: 1d-prefix-sum}을 보라.

\begin{code} \label{code: 1d-prefix-sum}
일차원 부분 합
\begin{lstlisting}
#include <bits/stdc++.h>
using namespace std;

int n, arr[101010], sum[101010], q;

int main() {
    cin >> n;
    for (int i = 1; i <= n; i++) cin >> arr[i];
    for (int i = 1; i <= n; i++) sum[i] = sum[i - 1] + arr[i];

    cin >> q;
    while (q--) {
        int s, e;
        cin >> s >> e;
        cout << sum[e] - sum[s - 1] << "\n";
    }

    return 0;
}
\end{lstlisting}
\end{code}

\section{이차원 부분 합}
\Cref{ex: 1d-sum-query}을 확장한 다음 문제를 생각해보자.

\begin{example}
    (11660번) 크기 $N \times M$인 어떤 격자 $A$에서 다음 쿼리 $Q$개를 해결하는 프로그램을 작성하라.
    \begin{itemize}
        \item 쿼리: $\sum_{i=s}^e \sum_{j=l}^r A_{i, j}$을 구하여라.
    \end{itemize}
\end{example}

이 문제에도 똑같이 부분 합의 원리를 적용할 수 있다. $S_{i, j}$를 $(1, 1)$부터 $(i, j)$까지의 합이라고 정의하자. 그러면 $(s, l)$부터 $(e, r)$까지의 합은 다음 식으로 구할 수 있다.

$$S_{e, r} - S_{s-1, r} - S_{e, l-1} + S_{s-1, l-1}$$

위 식이 성립한다는 것은 그림을 그리면 쉽게 알 수 있을 것이다. 첫 번째 $S_{e, r}$에 해당하는 영역을 칠하고 $S_{s-1, r}$ 영역에 해당하는 영역과 $S_{e, l-1}$에 해당하는 영역을 지워보라. 그러면 $S_{s-1, l-1}$에 해당하는 영역이 두 번 지워지므로 다시 더해주자. 그러면 원하는 영역이 된다.

마지막 남은 문제는 $S_{i, j}$를 구하는 것이다. 가장 쉬운 방법은 각 가로줄마다 부분 합을 구하고, 만들어진 새로운 배열에서 각 세로줄에 대해 부분 합을 구하는 것이다. 이 과정을 통해 만들어진 배열의 $(i, j)$번 칸은 $(1, 1)$부터 $(i, j)$까지의 합이 된다. 구현은 다음 코드 \Cref{code: 2d-prefix-sum}를 보라.

\begin{code} \label{code: 2d-prefix-sum}
이차원 부분 합
\begin{lstlisting}
#include <bits/stdc++.h>
using namespace std;

int n, arr[1010][1010], sum[1010][1010], q;

int main() {
    cin >> n;
    for (int i = 1; i <= n; i++) {
        for (int j = 1; j <= m; j++) {
            cin >> arr[i][j];
        }
    }
    for (int i = 1; i <= n; i++) {
        for (int j = 1; j <= m; j++) {
            sum[i][j] = sum[i][j - 1] + arr[i][j];
        }
    }
    for (int i = 1; i <= n; i++) {
        for (int j = 1; j <= m; j++) {
            sum[i][j] = sum[i - 1][j] + sum[i][j];
        }
    }

    cin >> q;
    while (q--) {
        int s, e, l, r;
        cin >> s >> l >> e >> r;
        cout << sum[e][r] - sum[s - 1][r] - sum[e][l - 1] + sum[s - 1][l - 1] << "\n";
    }

    return 0;
}
\end{lstlisting}
\end{code}

\section{시간복잡도 줄이기}

부분 합이 중요한 이유는 시간복잡도를 획기적으로 줄일 수 있기 때문이다. 다음 문제를 보자.

\begin{example}
    (13586번) `Go--' 게임은 $N \times N$의 바둑판 위에 흑과 백 모두 $P$개의 돌을 놓은 후, 바둑판에 있는 모든 부분 \textbf{정사각형}들 중 자기 색의 돌을 1개 이상 포함하며 상대 색의 돌을 포함하지 않는 정사각형의 개수가 더 많은 사람이 이기는 간단한 바둑 변형 게임이다. 바둑판의 정보가 주어질 때 흑과 백에 대해 정사각형의 개수를 계산해 출력하라. $N \leq 500$, $P \leq \min(\frac{N}{2}, 500)$이다.
\end{example}

부분 합을 사용하지 않으면, 고려해야 하는 정사각형의 수는 $O(N^3)$이고 각 돌의 개수를 세는데 $O(N^2)$의 시간이 필요하므로 $O(N^5)$이 된다. 하지만, 부분합을 사용하면 $O(1)$의 각 돌의 개수를 셀 수 있어 $O(N^3)$에 답을 구할 수 있다.

\section{누적 합의 역연산}

다음 문제를 보자.

\begin{example}
    (10713번) $i$번 철도역과 $i+1$번 철도역 사이를 잇는 $i$번 노선은 IC카드 구매 없이 $A_i$원에 이용하거나 IC카드 구매 후 $A_i$ 원보다 싼 가격인 $B_i$원에 이용할 수 있다. IC카드의 구매 비용은 $C_i$이며 IC카드는 노선마다 고유하다. 여러분은 여행 시작 전 일부 노선의 IC카드를 구매해 여행에 드는 교통비를 최소화하고 싶다. $M$일 동안 여행할 도시의 목록 $P_j$가 주어질 때, 여행 시작 전 IC카드를 사는 비용을 포함하여 필요한 교통비의 최솟값을 출력하라. $2 \leq N \leq 10^5$, $2 \leq M \leq 10^5$, $1 \leq B_i < A_i \leq 10^5$, $1 \leq C_i \leq 10^5$이다.
\end{example}

때로는 누적 합을 통해 원하는 배열을 얻어야 하는 경우가 있다. 다음 관찰을 생각하자.

\begin{obs}
    $A_0 = 0$이라 하고, $B_i = A_i - A_{i-1}$로 정의한 배열에 대해, $\sum_{i=1}^k B_i = A_k~(k \geq 1)$가 성립한다.
\end{obs}

결국 얻어야 하는 배열 $A_i$에 대해, $i$번째에 $A_i - A_{i-1}$을 저장하면 누적 합을 시행했을 때 $A_i$가 남게 된다. 이것을 활용하면 시간복잡도를 줄일 수 있다.

일단, IC카드를 구매하는 것이 유리할 조건은 $i$번 노선을 타는 횟수 $x$에 대해 $A_i \cdot x > B_i \cdot x + C_i$와 동치이다. 즉, $i$번 노선을 타는 횟수를 빠르게 구하면 문제를 해결할 수 있다. 이제 $i$번 노선을 타는 횟수를 생각해보면 $P_{j-1}$부터 $P_j$까지는 순차적으로 이동해야 하기 때문에, $[P_{j-1}, P_j)$ 또는 $[P_j, P_{j-1})$에 모두 $1$을 더해야 한다.

만약 직접 모든 칸에 $1$씩 더한다면 $O(NM)$의 시간이 필요하다. 하지만 연속된 구간에 전부 같은 수를 더하기 때문에, $B_i = A_i - A_{i-1}$에 해당하는 차이가 바뀌는 지점은 양 끝 두 지점이 전부이다. 따라서 $B_i$는 $O(M)$의 시간에 계산할 수 있다. 즉, $A_i$를 $O(N+M)$의 시간에 구할 수 있다.

\section{연습문제 A}
\begin{problem}
    \Cref{code: 1d-prefix-sum}의 시간복잡도는 무엇인가?
\end{problem}
\begin{problem}
    \Cref{ex: 1d-sum-query}에서 $A_i$를 $x$로 변경하는 새로운 종류의 쿼리가 추가된다면 문제를 어떻게 해결해야 할까?
    (이는 \Cref{part: 8}에 나오는 내용으로, 답을 찾지 못해도 상관없다. 궁금한 사람은 \Cref{part: 8}을 먼저 보라.)
\end{problem}
\begin{problem}
    \Cref{code: 2d-prefix-sum}의 시간복잡도는 무엇인가?
\end{problem}
\begin{problem}
    삼차원 부분 합은 식이 어떻게 될까? $d$차원에서는?
\end{problem}

\section{연습문제 B}

\begin{problem}
    (8694번*) $N$개의 물잔에 $w_i$만큼의 물이 들어 있다. \textbf{인접한} 물컵에 물을 옮기는 것만으로 모든 물컵에 든 물의 양을 똑같이 하려고 할 때 필요한 최소 물 이동 횟수를 출력하라. 물컵의 최대 크기는 무한하다고 가정할 수 있고 $w_i$의 합은 $N$의 배수이다. $N \leq 10^6$이고 $w_i$는 $10^6$인 정수이다.
\end{problem}

\begin{problem}
    (11525번) 분모가 $N$ 이하의 자연수이고 크기가 $0$ 이상 $1$ 이하인 기약분수를 크기 순으로 나열한 수열을 $N$번째 페리 수열이라고 한다. $N$번째 페리 수열의 길이를 구하는 $Q$개의 쿼리를 처리하는 프로그램을 작성하라. $N \leq 10000$, $Q \leq 10000$이다. (Hint. 오일러 토션트 함수)
\end{problem}

\begin{problem}
    (18628번*) 각 칸에 알파벳 소문자가 하나 씩 적혀 있는 $N \times M$ 격자의 좌상단에서 우하단까지 오른쪽 또는 아래로만 이동하여 도달하는 서로 다른 경로 세 개가 경로상에 있는 알파벳을 나열한 문자열이 같도록 잘 고를 수 있는지 판별하라. $N \leq 3000$, $M \leq 3000$이다.
\end{problem}

\begin{problem}
    (22046번, IOI21) \texttt{A, T, C}로만 이루어진 두 DNA 염기서열 $A$와 $B$ 사이의 돌연변이 거리는 $A$의 염기서열의 두 문자를 바꾸는 작업만으로 염기서열 $B$를 만들기 위해 필요한 최소 작업 횟수로 정의된다. 여러분은 길이 $N$의 염기서열 $A$와 $B$가 주어질 때 $Q$번 $(x, y)$를 입력받아 부분 염기서열 $A[x..y]$와 $B[x..y]$ 사이의 돌연변이 거리를 구해야 한다.
    \begin{itemize}
        \item (서브테스크 1) $N \leq 10^5$, $Q \leq 10^5$, $y - x \leq 2$
        \item (서브테스크 2) $N \leq 10^5$, $Q \leq 500$, $y - x \leq 1000$, $A$와 $B$를 이루는 문자는 \texttt{A, T} 뿐이다.
        \item (서브테스크 3) $N \leq 10^5$, $Q \leq 10^5$, $A$와 $B$를 이루는 문자는 \texttt{A, T} 뿐이다.
        \item (서브테스크 4) $N \leq 10^5$, $Q \leq 500$, $y - x \leq 1000$
        \item (서브테스크 5) $N \leq 10^5$, $Q \leq 10^5$
    \end{itemize}
\end{problem}

\begin{problem}
    (28167번*) $N+1 \times M+1$ 크기의 격자가 있다. $1 \leq i \leq N$과 $1 \leq j \leq M$에서 각 $i$번 행에 대해 왼쪽 $j$칸을 칠할 확률 $p_{ij}$가 주어지고, 각 $j$번 열에 대해 위쪽 $i$칸을 칠할 확률 $q_{ij}$가 주어진다. 그리고 $\sum_j p_{ij}$와 $\sum_i q_{ij}$는 $1$이다.
    색으로 구분된 구역의 총 개수의 기댓값을 구하라. 같은 구역은 같은 색의 칸이 상하좌우로 인접하게 연결되어 있는 것에 한한다. $N \leq 1000$, $M \leq 1000$이고 모든 확률은 유리수에 대한 $\textrm{mod}~998~244~353$ 표현으로 주어지며, 이 형식으로 출력해야 한다.
\end{problem}

\section{연습문제 C}
\begin{problem}
    (1792번) 다음 쿼리 $N$개를 처리하라. $N \leq 50000$이다.
    \begin{itemize}
        \item 세 정수 $a, b, d$가 주어질 때, $1 \leq x \leq a$, $1 \leq y \leq b$, $\gcd(x, y)=d$인 순서쌍 $(x, y)$의 개수를 구하여라.
        \item $1 \leq d \leq a \leq 50000$, $1 \leq d \leq b \leq 50000$
    \end{itemize}
\end{problem}

\begin{problem}
    $N$개의 방이 있는 호텔에서 매 손님이 올 때마다 랜덤으로 비어 있는 인접한 두 개의 방을 할당한다. 이때 가능한 각 인접한 방 쌍이 할당될 확률은 모두 같다. 더 이상 빈 인접한 두 방이 존재하지 않을 때 $k$번 방이 비어있을 확률을 $\textrm{mod}~10^9+7$에 대한 표현으로 출력하라. 총 $T$개의 케이스에 대해 $(N, k)$ 쌍을 입력받아 답을 출력해야 한다.
    \begin{itemize}
        \item (14340번) $T \leq 100$, $N \leq 10^4$
        \item (14341번) $T \leq 100$, $N \leq 10^7$
    \end{itemize}
\end{problem}

\begin{problem}
    (16001번) 길이 $N$인 문자열이 주어질 때, 몇 개의 문자를 바꾸어 역순 문자열이 정순 문자열과 같은 팰린드롬을 만드려고 한다. 이때 비용은 다음과 같이 계산된다.
    \begin{itemize}
        \item \textbf{현재 포인터 $p$가 $i$번째 문자를 가리킨다면} $i$번째 문자를 $A_i$의 비용으로 원하는 알파벳으로 교체할 수 있다.
        \item 현재 포인터 $p$가 $i$번째 문자를 기리킬 때, $j$번째 문자를 가리키게 하고 싶다면 $c \times \vert{i-j}\vert$의 비용이 필요하다.
    \end{itemize}
    처음 포인터 $p$가 가리키는 문자가 $k$번째 문자일 때 팰린드롬을 만들기 위해 필요한 최소 비용을 모든 $1 \leq k \leq N$에 대해 계산하라. $N \leq 10^6$, $A_i$와 $c$는 $10^9$ 이하의 자연수이다.
\end{problem}

\begin{problem}
    (19509번**) $T$번 $N$이 주어질 때마다 다음을 만족하는 길이 $N$인 수열을 하나 출력하라. $T \leq 200$, $N \leq 2000$이다.
    \begin{itemize}
        \item 수열의 모든 값은 $3(N+6)$ 이하의 양의 정수여야 한다.
        \item $1 \leq l \leq r \leq N$을 만족하는 모든 $N(N+1)/2$개의 정수 순서쌍 $(l, r)$에 대해 $l$번째 수부터 $r$번째 수까지의 합을 구하면 모든 순서쌍에서 전부 다르다.
    \end{itemize}
\end{problem}

\begin{problem}
    (20581번) 길이 $N$의 바이너리 문자열(\texttt{0, 1}로만 이루어진 문자열)은 총 $2^N$개 있다. 길이 $N$의 세 바이너리 문자열 $S_1, S_2, S_3$와 음 아닌 정수 $r_1, r_2, r_3$가 입력으로 들어올 때 다음 조건을 만족하는 길이 $N$의 바이너리 문자열의 개수를 출력하는 프로그램을 작성하라. $N \leq 10000$이다.
    \begin{itemize}
        \item 두 바이너리 문자열 $S, T$의 \textit{바이너리 거리}를 $\sum_i {\vert S_i - T_i \vert}$로 정의할 때, $S_i$과의 바이너리 거리가 $r_i$ 이하인 $i$가 존재해야 한다.
    \end{itemize}
\end{problem}

\newpage

\section{힌트}

\begin{hint}
    
\end{hint}

\begin{hint}
    
\end{hint}

\begin{hint}
    
\end{hint}

\begin{hint}
    
\end{hint}

\begin{hint}
    왼쪽부터 바로 오른쪽으로 물을 넘기되, 음수량을 넘겨도 상관이 없다. 순서를 조금 바꿔서 오른쪽에 물이 충분히 차면 그때 반대로 넘기면 되기 때문이다. 즉, 답의 상한은 $N-1$이다.
\end{hint}

\begin{hint}
    $1$ 이상 $n$ 이하의 자연수 중에서 $n$과 서로소인 수의 개수를 나타내는 함수를 오일러 토션트 함수라 하고, $\phi(n)$이라 쓴다. 이를 이용하면 분모가 $N$ 이하인 페리 수열의 길이를 나타낼 수 있다. 참고로,
    $$\textrm{factorization:}~n = p_1^{a_1} p_2^{a_2} \cdots p_k^{a_k} \Longrightarrow \phi(n) = n \prod_{i=1}^k \left( 1 - \frac{1}{p_i} \right)$$
    가 성립한다.
\end{hint}

\begin{hint}
    두 개의 단계별 힌트가 있다.

    \textbf{[ 1 ]} \newline
    같은 알파벳이 나열된 다른 경로가 나오려면, 어떤 칸을 기준으로 대각선 왼쪽 아래에 같은 알파벳이 있어야 한다. 따라서, 어떤 칸에 대해 바로 위 칸에 적힌 알파벳과 바로 왼쪽 칸에 적힌 알파벳이 같은 칸을 먼저 확인하자.

    \textbf{[ 2 ]} \newline
    \textbf{[ 1 ]}에서 확인한 칸을 \textit{좋은 칸}이라 하자. 그림을 그려보면 적어도 좋은 칸이 두 개 이상 존재해야 하고, 각각의 위치를 $(r_1, c_1)$과 $(r_2, c_2)$라 할 때, $r_1 < r_2, c_1 < c_2$ 또는 $(r_1, c_1) = (r_2, c_2-1)$ 또는 $(r_1, c_1) = (r_2-1, c_2)$이어야 한다.
\end{hint}

\begin{hint}
    주어진 구간에서 \texttt{A, C, T}의 개수가 다르면 불가능하다. 따라서 이 경우에는 \texttt{-1}을 리턴하면 된다. 이제 답이 있는 경우만을 생각하자.

    어떤 두 문자에 대해 상호 위치가 변경된 경우에는 한 번의 시행으로 해결할 수 있다. 이외의 경우는 \texttt{A, C, T}가 순환하는 경우이므로 한 쌍당 두 번의 시행으로 해결할 수 있다. 이것이 최적이므로 누적합을 이용해 상호 위치 변경 횟수, 순환하는 개수를 세면 쿼리당 $O(1)$에 해결할 수 있다.
\end{hint}

\begin{hint}
    총 두 개의 단계별 힌트가 있다.
    
    \textbf{[ 1 ]} \newline
    색칠된 같은 구역은 언제나 하나이다. 각 가로줄에서 왼쪽에서부터 적어도 하나는 칠하고, 각 세로줄에서 위쪽에서부터 적어도 하나는 칠하기 때문이다.

    \textbf{[ 2 ]} \newline
    어떤 칸이 색칠되어 있다면, 적어도 위쪽 칸이나 왼쪽 칸 중 하나는 색칠되어 있음이 보장된다. 따라서, 색칠되지 않은 같은 구역과, 자신은 색칠되지 않았으면서 아래 칸과 오른쪽 칸이 모두 색칠된 칸은, 일대일대응된다. 이러한 칸의 수는 누적 합을 사용하면 $O(NM)$에 빠르게 계산할 수 있다.
\end{hint}

\begin{hint}
    총 세 개의 단계별 힌트가 있다.

    \textbf{[ 1 ]} \newline
    뫼비우스 반전 공식을 사용해야 한다. 뫼비우스 반전 정리에 의하면

    $$g(n) = \sum_{d|n} f(d) \Rightarrow f(n) = \sum_{d|n} \mu(d) g\left( \frac{n}{d} \right)$$

    이 성립하는데, $\epsilon(n) = \begin{cases} 1 & (n=1) \\ 0 & (n>1) \end{cases}$과 $1(n) = 1$에 대해 $1(n) = \sum_{d|n} \epsilon(d)$이므로 $\epsilon(n) = \sum_{d|n} \mu(d) 1\left( \frac{n}{d} \right)$이다.

    \textbf{[ 2 ]} \newline
    구하는 값은 
    $$\sum_{i=1}^{a/d} \sum_{j=1}^{b/d} \epsilon(\gcd(i, j))$$
    이고, 뫼비우스 반전 공식에 의해
    $$\sum_{i=1}^{a/d} \sum_{j=1}^{b/d} \sum_{e|gcd(i, j)} \mu(e)$$
    와 같고, $e$를 기준으로 정리하면 $i$와 $j$가 모두 $e$의 배수인 경우 $\mu(e)$가 더해질 것이므로
    $$\sum_{e \geq 1} \left\lfloor \frac{a}{de} \right\rfloor \left\lfloor \frac{b}{de} \right\rfloor \mu(e)$$
    과 같다.

    \textbf{[ 3 ]} \newline
    \textbf{[ 2 ]}의 식을 그냥 계산한다면 $O(a/d)$의 시간이 소요되므로, $O(N \times a/d)$의 시간복잡도가 되어 TLE를 받는다. 여기서, $\left\lfloor \frac{a}{de} \right\rfloor$의 값의 종류의 수는 적음을 이용할 수 있다. 예를 들어 $\left\lfloor \frac{10}{e} \right\rfloor$을 $e$에 따라 적으면 $10, 5, 3, 2, 2, 1, 1, 1, 1, 1$로 다섯 종류의 수만 나온다. 수의 종류의 수는 $O(\sqrt{a/d})$이므로(harmonic lemma) 같은 수를 묶어서 계산하면 $O(\sqrt{a/d})$의 시간에 \textbf{[ 2 ]}의 식을 계산할 수 있다.
\end{hint}

\begin{hint}
    \url{https://github.com/kamyu104/GoogleCodeJam-2016/blob/master/World%20Finals/family-hotel.py}
\end{hint}

\begin{hint}
    \url{https://tech.kakao.com/posts/354}
\end{hint}

\begin{hint}
    인용: \url{https://m.blog.naver.com/PostView.naver?blogId=vermeil_&logNo=223146073298&navType=by}

    2년 간의 싸움 끝에 승리했다.

    x = (N보다 큰 소수)로 정하자. 그리고 누적 합 배열을 만드는데,

    pf[i] = i * p * 2 + (i * (i + 1) / 2) * x % p

    로 정의하자. 이는 우리가 만든 수열의 1~i번째 값의 합이다. 그리고 놀랍게도 저렇게 만들면 합이 안 겹친다 ㅋㅋ;

    pf[x] - pf[y] = pf[z] - pf[w] 를 풀면, 해는 (x = z) and (y = w)가 된다.

    이제 구현만 열심히 하면 끝

    \url{https://github.com/jbkarvens/baekjoon/tree/main/%EB%B0%B1%EC%A4%80}
\end{hint}

\begin{hint}
    \url{https://qoj.ac/contest/501/problem/1245/statistics}
\end{hint}